\chapter*{Prakata}

% SOURCE-CORRECTION: PREFACE-C001
Dalam karyanya yang elok, \emph{Hilbert Space Problem
Book}~\cite{Halmos:1982}, Paul Halmos terkenal menekankan bahwa praktik aktif merupakan
syarat mutlak untuk belajar matematika. Halmos tentu bukan satu-satunya yang meyakini hal
itu. Kumpulan catatan ini merupakan pendamping berorientasi aktivitas
untuk mempelajari analisis fungsional linear dan aljabar operator. Catatan ini dimaksudkan
sebagai pendamping pedagogis bagi pemula, pengantar menuju beberapa gagasan utama dalam
bidang analisis ini, kumpulan soal yang menurut saya berguna untuk mempelajari pokok
bahasannya, serta daftar bacaan dan referensi beranotasi.

Sebagian besar hasil dalam pengantar analisis fungsional bersifat lugas dan dapat
diperiksa oleh mahasiswa yang menelaahnya dengan saksama. Bahkan, itulah tujuan utama
catatan ini---meyakinkan pemula bahwa bidang ini dapat dipelajari. Dalam materi berikut
terdapat banyak penanda yang mengajak pembaca untuk berperan aktif: kata-kata seperti
``proposisi'', ``contoh'', ``latihan'', dan ``akibat'', jika tidak diikuti oleh bukti
atau rujukan kepada suatu bukti, merupakan undangan untuk memeriksa pernyataan yang
diajukan. Tentu saja, ada beberapa teorema yang, menurut saya, memerlukan waktu dan usaha
pembuktian yang lebih besar daripada manfaat yang diperoleh. Dalam kasus seperti itu,
apabila saya tidak dapat menyempurnakan bukti bakunya, saya memberikan rujukan alih-alih
mengulang bukti tersebut.

Materi dalam catatan ini mencerminkan isi sebagian dari dua mata kuliah pascasarjana
sepanjang tahun tentang beberapa topik dasar dalam \emph{analisis fungsional (linear)}
dan \emph{aljabar operator}, yang telah berkali-kali saya ajarkan di Portland State
University. Dalam mata kuliah tersebut, selain menggunakan catatan ini, saya meminta
mahasiswa memilih sumber informasi lain---baik buku cetak maupun dokumen daring---yang
sesuai dengan gaya belajar (dan anggaran) masing-masing. Prasyarat untuk mempelajari
catatan ini cukup ringan; bahkan pengenalan sepintas terhadap \emph{aljabar linear},
\emph{analisis real}, dan \emph{topologi} semestinya memadai. Dengan kata lain, istilah
seperti \emph{ruang vektor}, \emph{pemetaan linear}, \emph{limit}, \emph{ukuran Lebesgue}
dan \emph{integral Lebesgue}, \emph{himpunan terbuka}, \emph{himpunan kompak}, serta
\emph{fungsi kontinu} semestinya sudah tidak asing.

Karena \emph{analisis fungsional linear}, setidaknya dalam suatu pengertian, dapat
dipandang sebagai `aljabar linear berdimensi tak hingga', bab pertama catatan ini
dikhususkan bagi tinjauan (yang cukup ringkas) atas beberapa gagasan penting yang,
  dalam dunia angan-angan saya, dibawa pulang mahasiswa dari mata kuliah
aljabar linear pertama mereka. Tiga wawasan utama yang dikembangkan dalam bab pertama
adalah:
\begin{enumerate}
  \item operator yang dapat didiagonalkan pada ruang vektor berdimensi hingga merupakan
  kombinasi linear dari proyeksi;
  \item operator pada ruang hasil kali dalam kompleks berdimensi hingga dapat didiagonalkan secara
  uniter jika dan hanya jika operator itu normal; dan
  \item dua operator normal pada ruang hasil kali dalam kompleks berdimensi hingga ekuivalen secara
  uniter jika dan hanya jika keduanya memiliki nilai-nilai eigen yang sama, masing-masing
  dengan multiplisitas yang sama. % SOURCE-CORRECTION: PREFACE-C002
\end{enumerate}
Saya memandang hasil-hasil tersebut sangat mendasar karena wawasan luar biasa yang lahir
dari berbagai upaya para ahli teori operator selama seabad terakhir untuk membawanya ke
dalam ranah \emph{analisis fungsional}, yakni menggeneralisasikannya ke kasus berdimensi
tak hingga. Salah satu akibat penjelajahan yang subur dan luar biasa rumit ini ialah
ditemukannya banyak hubungan yang mencengangkan, erat, dan kompleks antara \emph{teori
operator} dan bidang-bidang seperti \emph{teori fungsi variabel kompleks}, \emph{topologi
aljabar}, serta \emph{aljabar homologis}. Dalam bagian-bagian selanjutnya dari catatan ini, saya
berusaha menelusuri kembali beberapa langkah tersebut. Kuliah yang menjadi dasar catatan
ini berpuncak pada hasil menakjubkan Brown, Douglas, dan Fillmore yang pertama kali
diterbitkan pada 1973 (lihat~\cite{BrownDF:1973}). Bahkan sekarang, sebagian besar mahasiswa
pascasarjana masih kesulitan menemukan uraian yang cukup mudah dicerna mengenai banyak
materi latar yang diperlukan untuk mengembangkan materi indah ini. Sayangnya, penalaran
transfinit yang serba informal dan masih cukup berhasil di ruang kuliah sangat sukar dialihkan ke kertas
atau media elektronik. Akibatnya, catatan ini masih berada dalam keadaan yang sungguh
belum selesai. Saya terus mengerjakannya; dan mungkin orang lain kelak akan menyelesaikan
apa yang tidak saya selesaikan.

Jika materi aljabar linear dalam bab pertama terasa terlalu ringkas, uraian yang lebih
santai dan lebih menyeluruh dapat ditemukan dalam catatan daring saya~\cite{Erdman:2010}.
% SOURCE-CORRECTION: PREFACE-C003
Demikian pula, materi latar tambahan tentang kalkulus lanjut, ruang metrik dan ruang
topologis, integral Lebesgue, dan sebagainya dapat ditemukan dalam banyak sekali sumber,
termasuk~\cite{Erdman:2005} dan~\cite{Erdman:2007}.

Tentu terdapat sejumlah kelebihan dan kekurangan ketika sebuah dokumen diterbitkan secara
elektronik. Salah satu kelebihannya ialah rujuk silang dapat diikuti dengan cepat melalui
pranala. Menelusuri sebuah buku untuk mencari \emph{lema 3.14.8} dapat memakan waktu
(terutama apabila seorang penulis menerapkan kebiasaan yang sepenuhnya logis tetapi
sangat menjengkelkan, yakni memberi nomor secara terpisah kepada lema, proposisi, teorema,
akibat, dan sebagainya). Kelebihan yang mungkin lebih penting ialah kemampuan untuk
segera memperbaiki kesalahan, menambahkan bagian yang hilang, memperjelas argumen yang
samar, dan membenahi ungkapan yang kurang baik. Kekurangan yang menyertainya ialah bahwa
seorang pembaca yang kembali ke laman web setelah waktu singkat mungkin mendapati segala
sesuatu (halaman, definisi, teorema, bagian) telah memperoleh nomor yang berbeda.
({\LaTeX} adalah alat yang mengagumkan.) Saya akan mengubah tanggal pada halaman judul
untuk memberi tahu pembaca kapan pembaruan tak sepele terakhir dilakukan (yakni pembaruan
yang memengaruhi nomor atau rujuk silang).

Kekurangan paling serius dari kehidupan elektronik ialah ketidaklestariannya. Dalam
kebanyakan kasus, ketika sebuah laman web menghilang, untuk segala keperluan praktis
informasi di dalamnya turut menghilang. Karena alasan ini (dan karena saya ingin materi
ini tersedia secara bebas bagi siapa pun yang menginginkannya), saya menggunakan lisensi
``BerbagiSerupa'' dari \emph{Creative Commons}. % SOURCE-CORRECTION: PREFACE-C004
Saya berharap siapa pun yang menganggap
materi ini berguna akan memperbaiki yang keliru, menambahkan yang hilang, dan menyempurnakan
yang canggung. Saya akan menyertakan berkas sumber {\LaTeX} di situs web saya agar
pemanfaatan dokumen ini, atau bagian-bagiannya, tidak mengharuskan pengetikan ulang tanpa
henti. Mengenai teksnya sendiri, silakan kirim koreksi, saran, keluhan, dan komentar lain
kepada penulis melalui
   \[ \textrm{erdman@pdx.edu} \]


\vfill\eject


\section*{Huruf Yunani}

 \index{huruf!Yunani}%
% SOURCE-CORRECTION: PREFACE-C005
% Markup tabel warisan yang dikecualikan diganti dengan tabularx; data dipertahankan.
\begin{center}
{\renewcommand{\arraystretch}{1.10}%
\begin{tabularx}{\textwidth}{@{}>{\centering\arraybackslash}p{0.23\textwidth}>{\centering\arraybackslash}p{0.27\textwidth}>{\raggedright\arraybackslash}X@{}}
\hline
\textbf{Huruf kapital} & \textbf{Huruf kecil} & \textbf{Nama dalam bahasa Inggris (perkiraan pelafalan)} \\
\hline
A         & $\alpha$                    & Alpha (AL-fuh) \\
B         & $\beta$                     & Beta (BAY-tuh) \\
$\Gamma$  & $\gamma$                    & Gamma (GAM-uh) \\
$\Delta$  & $\delta$                    & Delta (DEL-tuh) \\
$E$       & $\epsilon$ atau $\varepsilon$ & Epsilon (EPP-suh-lon) \\
Z         & $\zeta$                     & Zeta (ZAY-tuh) \\
H         & $\eta$                      & Eta (AY-tuh) \\
$\Theta$  & $\theta$                    & Theta (THAY-tuh) \\
I         & $\iota$                     & Iota (eye-OH-tuh) \\
K         & $\kappa$                    & Kappa (KAP-uh) \\
$\Lambda$ & $\lambda$                   & Lambda (LAM-duh) \\
M         & $\mu$                       & Mu (MYOO) \\
N         & $\nu$                       & Nu (NOO) \\
$\Xi$     & $\xi$                       & Xi (KSEE) \\
O         & o                            & Omicron (OHM-ih-kron) \\
$\Pi$     & $\pi$                       & Pi (PIE) \\
P         & $\rho$                      & Rho (ROH) \\
$\Sigma$  & $\sigma$                    & Sigma (SIG-muh) \\
T         & $\tau$                      & Tau (TAU) \\
Y         & $\upsilon$                  & Upsilon (OOP-suh-lon) \\
$\Phi$    & $\phi$                      & Phi (FEE or FAHY) \\
X         & $\chi$                      & Chi (KHAY) \\
$\Psi$    & $\psi$                      & Psi (PSEE or PSAHY) \\
$\Omega$  & $\omega$                    & Omega (oh-MAY-guh) \\
\hline
\end{tabularx}}
\end{center}








\vfill\eject







\section*{Font Fraktur}

Dalam catatan ini digunakan font Fraktur (paling sering untuk keluarga himpunan dan
keluarga operator). Padanan huruf Roman bagi setiap huruf diberikan di bawah ini. Ketika
menulis dengan tangan atau menyajikan materi di papan tulis, biasanya paling baik
menggantinya dengan huruf Latin kaligrafis.
 \index{font!Fraktur}

% SOURCE-CORRECTION: PREFACE-C006
% Markup tabel warisan yang dikecualikan diganti dengan tabularx; data dipertahankan.
\begin{center}
{\renewcommand{\arraystretch}{1.10}%
\begin{tabularx}{\textwidth}{@{}>{\centering\arraybackslash}X>{\centering\arraybackslash}X>{\centering\arraybackslash}X@{}}
\hline
\textbf{Fraktur} & \textbf{Fraktur} & \textbf{Roman} \\
\textbf{Huruf kapital} & \textbf{Huruf kecil} & \textbf{Huruf kecil} \\
\hline
$\mathfrak A$ & $\mathfrak a$ & a \\
$\mathfrak B$ & $\mathfrak b$ & b \\
$\mathfrak C$ & $\mathfrak c$ & c \\
$\mathfrak D$ & $\mathfrak d$ & d \\
$\mathfrak E$ & $\mathfrak e$ & e \\
$\mathfrak F$ & $\mathfrak f$ & f \\
$\mathfrak G$ & $\mathfrak g$ & g \\
$\mathfrak H$ & $\mathfrak h$ & h \\
$\mathfrak I$ & $\mathfrak i$ & i \\
$\mathfrak J$ & $\mathfrak j$ & j \\
$\mathfrak K$ & $\mathfrak k$ & k \\
$\mathfrak L$ & $\mathfrak l$ & l \\
$\mathfrak M$ & $\mathfrak m$ & m \\
$\mathfrak N$ & $\mathfrak n$ & n \\
$\mathfrak O$ & $\mathfrak o$ & o \\
$\mathfrak P$ & $\mathfrak p$ & p \\
$\mathfrak Q$ & $\mathfrak q$ & q \\
$\mathfrak R$ & $\mathfrak r$ & r \\
$\mathfrak S$ & $\mathfrak s$ & s \\
$\mathfrak T$ & $\mathfrak t$ & t \\
$\mathfrak U$ & $\mathfrak u$ & u \\
$\mathfrak V$ & $\mathfrak v$ & v \\
$\mathfrak W$ & $\mathfrak w$ & w \\
$\mathfrak X$ & $\mathfrak x$ & x \\
$\mathfrak Y$ & $\mathfrak y$ & y \\
$\mathfrak Z$ & $\mathfrak z$ & z \\
\hline
\end{tabularx}}
\end{center}


  \vfill\eject
\phantomsection
\section*{Notasi untuk Himpunan Bilangan}\label{C0009} % SOURCE-CORRECTION: PREFACE-C007
Berikut adalah daftar notasi yang cukup baku untuk beberapa himpunan bilangan yang sering
muncul dalam catatan ini:
 \index{real!garis!subhimpunan khusus dari}%
 \index{bilangan!himpunan khusus dari}%
 \index{C@$\C$ (himpunan bilangan kompleks)}%
 \index{R@$\R$ (himpunan bilangan real)}%
 \index{R@$\R^n$!himpunan tupel-$n$ bilangan real}%
 \index{Q@$\Q$ (himpunan bilangan rasional)}%
 \index{Z@$\Z$ (himpunan bilangan bulat)}%
 \index{N@$\N$ (himpunan bilangan asli)}%
 \index{N@$\N_n$ ($n$ bilangan asli pertama)}%
 \index{D@$\Di$ (cakram satuan terbuka)}%
 \index{T@$\T$ (lingkaran satuan)}%
 \index{S@$\Sp^1$ (lingkaran satuan)}%
\begin{align*}
 &\C \text{ adalah himpunan bilangan kompleks} \\
 &\R \text{ adalah himpunan bilangan real} \\
 &\R^n \text{ adalah himpunan semua tupel-$n$ $(r_1,r_2,\dots,r_n)$ bilangan real} \\
 &\R^+ = \{x \in \R\colon x \ge 0\}, \text{ himpunan bilangan real tak negatif} \\ % SOURCE-CORRECTION: PREFACE-C008
 &\Q \text{ adalah himpunan bilangan rasional} \\
 &\Q^+ = \{x \in \Q\colon x \ge 0\}, \text{ himpunan bilangan rasional tak negatif} \\ % SOURCE-CORRECTION: PREFACE-C009
 &\Z \text{ adalah himpunan bilangan bulat} \\
 &\Z^+ = \{0,1,2,\dots\}, \text{ himpunan bilangan bulat tak negatif} \\ % SOURCE-CORRECTION: PREFACE-C010
 &\N = \{1,2,3,\dots\}\text{, himpunan bilangan asli} \\
 &\N_n = \{1,2,3,\dots,n\}\text{, $n$ bilangan asli pertama} \\ % SOURCE-CORRECTION: PREFACE-C011
 &[a,b] = \{x \in \R\colon a \le x \le b\} \\
 &[a,b) = \{x \in \R\colon a \le x < b\} \\
 &(a,b] = \{x \in \R\colon a < x \le b\} \\
 &(a,b) = \{x \in \R\colon a < x < b\} \\
 &[a,\infty) = \{x \in \R\colon a \le x\} \\
 &(a,\infty) = \{x \in \R\colon a < x\} \\
 &(-\infty,b] = \{x \in \R\colon x \le b\} \\
 &(-\infty,b) = \{x \in \R\colon  x < b\} \\
\index{cakram!satuan terbuka}%
 &\Di = \{(x,y) \in \R^2\colon x^2 + y^2 < 1\},\text{ cakram satuan terbuka} \\
 &\T = \Sp^1 = \{(x,y) \in \R^2\colon x^2 + y^2 = 1\},\text{ lingkaran satuan} \\
\end{align*}



\vfill\eject










\section*{Notasi untuk Fungsi}
\enlargethispage{2\baselineskip}
Fungsi sudah tidak asing sejak kalkulus dasar. Secara informal, suatu \emph{fungsi} terdiri atas sepasang
himpunan dan suatu ``aturan'' yang mengaitkan setiap anggota himpunan pertama (yaitu
\emph{domain}) dengan tepat satu anggota himpunan kedua (yaitu \emph{kodomain}). Walaupun
``definisi'' informal ini tentu memadai untuk sebagian besar keperluan dan jarang menimbulkan
kesalahpahaman, kadang-kadang perumusan yang lebih tepat tetap berguna. Perumusan itu diperoleh
dengan mendefinisikan fungsi melalui suatu jenis khusus relasi antara dua himpunan.

Suatu
 \index{fungsi}%
\df{fungsi} $f$ adalah tripel terurut $(S,T,G)$, dengan $S$ dan $T$ himpunan serta $G$
subhimpunan dari $S \times T$ yang memenuhi:
 \begin{enumerate}
   \item untuk setiap $s \in S$ terdapat $t \in T$ sedemikian sehingga $(s,t) \in
G$; dan
   \item jika $(s,t_1)$ dan $(s,t_2)$ termasuk dalam $G$, maka $t_1 = t_2$. % SOURCE-CORRECTION: PREFACE-C014
 \end{enumerate}
Dalam keadaan ini kita mengatakan bahwa $f$ adalah \emph{fungsi dari} $S$ \emph{ke} $T$ (atau bahwa
$f$ \emph{memetakan} $S$ \emph{ke} $T$) dan menulis $f\colon S \sto T$. Himpunan $S$ adalah
 \index{domain}%
\df{domain} (atau
 \index{masukan!ruang}%
 \index{ruang!masukan}%
\df{ruang masukan}) dari~$f$. Himpunan $T$ adalah
 \index{kodomain}%
\df{kodomain} (atau
 \index{sasaran!ruang}%
 \index{ruang!sasaran}%
\df{ruang sasaran}, atau
 \index{keluaran!ruang}%
 \index{ruang!keluaran}%
\df{ruang keluaran}) dari~$f$. Relasi $G$ adalah
 \index{graf}%
\df{graf} dari~$f$. Untuk menghindari penyebutan graf $G$ secara eksplisit, biasanya
ungkapan ``$(x,y) \in G\,$'' diganti dengan ``$y = f(x)$''; elemen $f(x)$ adalah
 \index{citra!dari suatu titik}%
\df{citra} $x$ di bawah~$f$. Dalam catatan ini (tetapi tidak di semua tempat), kata
 \index{transformasi}%
 \index{konvensi!fungsi = peta = pemetaan = transformasi}%
``transformasi'',
 \index{peta}%
``peta'', dan
 \index{pemetaan}%
``pemetaan'' digunakan sebagai sinonim ``fungsi''. Domain $f$ dinotasikan
 \index{domain@$\dom f$ (domain fungsi $f$)}%
dengan~$\dom f$.

Jika $S$ dan $T$ himpunan, keluarga semua fungsi yang berdomain $S$ dan
berkodomain $T$ dinotasikan
 \index{F@$\fml F(S,T)$ (himpunan fungsi dari $S$ ke~$T$)}%
dengan~$\fml F(S,T)$.

Jika $f \colon S \sto T$ dan $A\subseteq S$, maka
 %\index{<@$\bigl.f\bigr"|_A$ (pembatasan $f$ pada $A$)}%
 \index{<@$f\mid_{{}_A}$ (pembatasan $f$ pada $A$)}%
 \index{pembatasan}%
\df{pembatasan} $f$ pada $A$, yang dinotasikan dengan $\bigl.f\bigr|_A$, adalah
pemetaan dari $A$ ke $T$ yang nilainya pada setiap $x$ di $A$ sama dengan $f(x)$.

Jika $f\colon S \sto T$ dan $A \subseteq S$, maka $f^{\sto} (A)$, yaitu
 \index{<arrowto@$f^{\sto}(A)$ (citra himpunan $A$ di bawah fungsi $f$)}%
 \index{citra!dari suatu himpunan}%
\df{citra} $A$ di bawah $f$ adalah $\{f(x)\colon x \in A\}$. Lazim pula menulis
$f(A)$ sebagai pengganti $f^{\sto}(A)$. Himpunan $f^{\sto}(S)$ adalah
 \index{jangkauan}%
 \index{jangkauan@$\ran f$ (jangkauan fungsi $f$)}%
\df{jangkauan} (atau
 \index{citra!dari suatu fungsi}%
\df{citra}) dari $f$; biasanya kita menulis $\ran f$ sebagai pengganti $f^{\sto}(S)$.

Jika $f\colon S \sto T$ dan $B \subseteq T$, maka
$f^\gets (B)$, yaitu
 \index{<arrowgets@$f^{\gets}(B)$ (pracitra himpunan $B$ di bawah fungsi $f$)}%
 \index{pracitra!dari suatu himpunan}%
 \index{citra!invers!dari suatu himpunan}%
\df{pracitra} $B$ di bawah~$f$ (atau citra inversnya) adalah $\{x \in S\colon f(x) \in B\}$. Dalam banyak buku,
$f^{\gets}(B)$ dinotasikan dengan $f^{-1}(B)$. Notasi ini dapat membingungkan karena menyiratkan
bahwa fungsi $f$ memiliki invers, padahal belum tentu demikian.

Suatu fungsi $f\colon S \sto T$ bersifat
 \index{injektif}%
\df{injektif} (atau
 \index{satu-satu}%
\df{satu-satu}) jika $f(x) = f(y)$ mengakibatkan $x = y$ untuk semua $x$, $y \in S$. Fungsi itu bersifat
 \index{surjektif}%
\df{surjektif} (atau
 \index{onto}%
\df{onto}) jika untuk setiap $t \in T$ terdapat $x \in S$ sedemikian sehingga $f(x) = t$. % SOURCE-CORRECTION: PREFACE-C012
Fungsi itu bersifat
 \index{bijektif}%
\df{bijektif} (atau merupakan
 \index{satu-satu!korespondensi}%
 \index{korespondensi!satu-satu}%
\df{korespondensi satu-satu}) jika fungsi itu injektif sekaligus surjektif. % SOURCE-CORRECTION: PREFACE-C013

Sering kali berguna untuk memandang fungsi sebagai panah dalam diagram. Sebagai contoh, keadaan
$j\colon R \sto U$, $f\colon R \sto S$, $k\colon S \sto T$, $h\colon U \sto T$
dapat direpresentasikan oleh
 \index{diagram}%
diagram berikut.
  \[ \xy
       \square[R`U`S`T;j`f`h`k]
  \endxy \]
Diagram itu dikatakan
 \index{berkomutasi}%
\df{berkomutasi} (atau disebut
 \index{komutatif!diagram}%
 \index{diagram!komutatif}%
\df{diagram komutatif}) jika $h \circ j = k \circ f$.

Diagram tidak harus berbentuk persegi panjang. Sebagai contoh,
 \[ \xy
     \btriangle[R`S`T;f`g`k]
 \endxy \]
adalah diagram komutatif jika $g = k \circ f$.


\vfill\eject













  \endinput
